\section{PR-II Exact Classifications and Canonical Closure}
\label{sec:pr2-branch-classification}

PR-II extends the PR-I foundation to an electroweak and flavor architecture \cite{PRII}.  The direct-ancestry audit separates exact consequences of a fixed electroweak representation from the logical closure principles used to construct that representation.

\subsection{Hypercharge Solution within Fixed Matter Content}

Write the one-generation left-handed Weyl hypercharges \cite{Weinberg_QFT_1996} as $q,u,d,\ell,e$ for the quark doublet, up conjugate, down conjugate, lepton doublet, and charged-lepton conjugate.  With $q=t$, the linear local-anomaly equations reduce to
\begin{equation}
 \ell=-3t,\qquad e=6t,\qquad d=-2t-u.
\end{equation}
Substitution into the cubic anomaly gives the exact factorization
\begin{equation}
 -18t(u-2t)(u+4t)=0.
 \label{eq:pr35-hypercharge-factorization}
\end{equation}
For $t\ne0$, the last two factors yield the standard branch and its up/down singlet relabeling.  Overall hypercharge normalization is conventional.  The $t=0$ factor gives an anomaly-free vectorlike singlet family $q=\ell=e=0$, $d=-u$; it is excluded only when the nonzero residual-charge ledger is imposed.  Hence anomaly cancellation alone is not a unique derivation, whereas anomaly cancellation plus a fixed nonzero charge ledger does isolate the familiar branch modulo normalization and labeling.

\subsection{Unique Photon Direction in the Declared Order Representation}

Once the neutral mass matrix is derived in outer-product form,
\begin{equation}
 M_0^2\propto
 \begin{pmatrix}g^2&-gg'\\-gg'&g'^2\end{pmatrix}
 =\begin{pmatrix}g\\-g'\end{pmatrix}
  \begin{pmatrix}g&-g'\end{pmatrix},
\end{equation}
nonzero $g,g'$ imply rank one and a one-dimensional nullspace.  The photon direction is unique up to normalization.  The outer-product form is a fixed inherited PR-II representation relation.  Other symmetric rank-one matrices violate that representation ledger and are outside the declared class.

\subsection{Generation Replication Under Anomaly Gates Alone}

Let one complete anomaly-free generation contain four weak doublets after color multiplicity.  Replicating the complete ledger $N$ times multiplies every local anomaly coefficient by $N$, so all remain zero.  The number of weak doublets is $4N$, which is even for every $N\in\mathbb Z_{>0}$; the
Witten \cite{Witten_1982} parity condition therefore also holds.  Consequently,
\begin{equation}
 N=1,2,3,\ldots
\end{equation}
is an infinite discrete family satisfying these gates.  Anomaly and Witten conditions alone therefore do not select $N_{\rm gen}=3$. Appendix~\ref{app:generation-count-rigidity} supplies the additional PR-II completeness/nonredundancy sheet-incidence theorem, exhausts every positive integer, and uniquely yields $N=3$ in the declared realization.  The normalized-incidence condition $\Phi_3^\dagger\Phi_3=I_3$ then leaves one unitary operator orbit, as proved in
Appendix~\ref{app:three-selector-closure}.  Its logical justification is minimal faithful realization: every broken compact direction must be represented (surjectivity), while no invisible sheet combination may remain (injectivity).  These are target-independent closure conditions for the
PR-II lift; literal appearance in P1 through P7 is not required.

The electroweak lift, weak boundary-deficit rule, hierarchy action, flavor overlaps, neutrino branch, and strong anchor likewise admit either direct ancestry proofs or explicit target-independent closure principles.  The graph records those two routes as complementary provenance classifications for the
PR-II extension.
